\section*{Problema 2}

Dado un conjunto de $n$-tuplas de vectores $A$ y $B$, se nos pide calcular el producto escalar $A \cdot B$ para cada una de las $n$-tuplas de los ejercicios 1 al 6 del \S 1.

El producto escalar de dos vectores $A = (a_1, a_2, \dots, a_n)$ y $B = (b_1, b_2, \dots, b_n)$ se define como:
\[
A \cdot B = \sum_{i=1}^n a_i b_i.
\]

A continuación, calculamos el producto escalar para cada uno de los pares de vectores dados:

\begin{enumerate}
    \item Para $A = (1, -1, 1)$ y $B = (2, 1, 5)$:
    \begin{align*}
        A \cdot B &= (1)(2) + (-1)(1) + (1)(5) \\
        &= 2 - 1 + 5 \\
        &= 6.
    \end{align*}

    \item Para $A = (1, -1, 1)$ y $B = (2, 3, 1)$:
    \begin{align*}
        A \cdot B &= (1)(2) + (-1)(3) + (1)(1) \\
        &= 2 - 3 + 1 \\
        &= 0.
    \end{align*}

    \item Para $A = (-5, 2, 7)$ y $B = (3, -1, 2)$:
    \begin{align*}
        A \cdot B &= (-5)(3) + (2)(-1) + (7)(2) \\
        &= -15 - 2 + 14 \\
        &= -3.
    \end{align*}

    \item Para $A = (\pi, 2, 1)$ y $B = (2, -\pi, 0)$:
    \begin{align*}
        A \cdot B &= (\pi)(2) + (2)(-\pi) + (1)(0) \\
        &= 2\pi - 2\pi + 0 \\
        &= 0.
    \end{align*}
\end{enumerate}